% IEEE BIBM-style conference manuscript draft.
% Main-paper target: IEEE two-column format, maximum eight pages including references.
% Revision v3, staged rewrite: this file currently contains the title and the
% rewritten abstract only. Remaining sections are carried over one at a time from
% the previous version (..._Risk_Controlled_Prediction_...tex).
% Main-paper figure assets are stored under figures/.
\documentclass[conference]{IEEEtran}
\IEEEoverridecommandlockouts

\usepackage{cite}
\usepackage{amsmath,amssymb,amsfonts}
\usepackage{algorithmic}
\usepackage{graphicx}
\usepackage{booktabs}
\usepackage{array}
\usepackage{textcomp}
\usepackage{xcolor}
\usepackage{url}
\usepackage{placeins}
\setlength{\emergencystretch}{2em}

% Compact float spacing for the two-column layout. These settings reduce
% whitespace around later figures without changing the manuscript text.
\setlength{\abovecaptionskip}{2pt}
\setlength{\belowcaptionskip}{0pt}
\setlength{\textfloatsep}{4pt plus 1pt minus 1pt}
\setlength{\floatsep}{3pt plus 1pt minus 1pt}
\setlength{\dbltextfloatsep}{4pt plus 1pt minus 1pt}
\setlength{\dblfloatsep}{3pt plus 1pt minus 1pt}

\def\BibTeX{{\rm B\kern-.05em{\sc i\kern-.025em b}\kern-.08em
    T\kern-.1667em\lower.7ex\hbox{E}\kern-.125emX}}

\newcommand{\placeholderfigure}[1]{%
  \fbox{\parbox[c][1.18in][c]{0.94\columnwidth}{\centering\footnotesize #1}}}
\newcommand{\placeholderfigurewide}[1]{%
  \fbox{\parbox[c][1.32in][c]{0.94\textwidth}{\centering\footnotesize #1}}}

\begin{document}

\title{Morphology-Evidence Factorized Diagnosis with Dual-Space Semantic Geometry and Risk-Aware Decision Support for Cervical Cytology}

\author{\IEEEauthorblockN{Wenjun Xu}}

\maketitle

\begin{abstract}
Errors in cervical cytology carry unequal risk: a high-grade lesion called
low-grade is a missed precursor, whereas ASC-US versus LSIL disagreement
reflects routine interpretive variation. A flat five-class classifier reports
uncertainty but not risk type, limiting differentiated review policies. This is
a representational constraint rather than an accuracy deficit, and higher
accuracy alone does not remove it. The Bethesda System defines atypical squamous
cell categories by insufficient evidence for definitive interpretation and
separates them by grade, implicitly factoring the four abnormal categories onto
a grade-related morphology axis and an evidence-definitiveness axis. We present
MERA-Dx (Morphology--Evidence Factorized and Risk-Aware Diagnostic Network),
which recovers these axes without additional annotation. A Swin-Tiny model is
trained with a direct five-class head and factor-specific supervision, fusing
beliefs only on the abnormal conditional distribution to preserve
Normal/Abnormal mass exactly. The frozen representation is then mapped into
separate morphology and evidence spaces with multi-prototype factor states.
Finally, a modular risk layer consumes frozen outputs to provide
screening-conflict detection, a factor--decision discordance sentinel, and
class-conditional pairwise prediction sets. In five-fold development evaluation,
the two axes separate as predicted: the evidence axis, only partially observable
from a single cell, shows a geometry gap of 0.5184 versus 1.2627 for morphology,
with factor accuracy near 0.76 versus 0.90. The sentinel refers 94.70\% of
hidden high-grade cases at a pre-specified operating point with 18.09\% routed
to safety-critical abstention within a 43.41\% review rate, and pairwise sets
reach approximately 95\% coverage. Both later layers leave the diagnosis
untouched, with zero parameter drift and zero prediction changes. Under matched
30-epoch training, Macro-F1 improves from 0.7675 to 0.7752. These are
development results; the risk layer is exploratory, and no external confirmation
cohort was available. The contribution is a factorization that turns an existing
taxonomy into an auditable basis for separating safety-critical from routine
uncertainty.
\end{abstract}

\begin{IEEEkeywords}
cervical cytology, Bethesda system, morphology-evidence factorization,
semantic geometry, prototype learning, conformal prediction,
clinical decision support, MERA-Dx
\end{IEEEkeywords}

\begin{figure}[t]
    \centering
    \includegraphics[width=\columnwidth]{figures/graphical_abstract.png}
    \caption{Conceptual overview. A flat five-class output can report that a case is uncertain but not which kind of uncertainty it carries. Reading the Bethesda System's own category definitions directly recovers a grade-related axis and an evidence-definitiveness axis, exposing a query a flat decision cannot pose and giving a basis for risk-differentiated review.}
    \label{fig:graphical_abstract}
\end{figure}

\section{Introduction}

Cervical cytology screening assigns every case to one of five Bethesda
System categories \cite{nayar2015bethesda}, and the clinical cost of an
error depends heavily on which two categories are confused. A high-grade lesion reported as
low-grade is a missed precursor that can delay treatment, while a
disagreement between two closely related atypical categories is closer
to ordinary interpretive variation and rarely changes clinical
management. A decision support system deployed into this workflow
inherits this asymmetry: it must not only be accurate, it must fail in
a way that a reviewing clinician can triage.

Automated classification of cervical cytology has largely been framed
as a single flat multi-class problem, where a backbone network outputs
one of the five category labels directly and progress is reported
through overall accuracy or macro-averaged F1. Uncertainty in this
framing is typically summarized by a single confidence or entropy score
attached to the predicted label, which answers whether the network is
unsure but not what kind of error it might be making. A flat classifier
has no internal variable that separates a grade-related mistake from an
evidence-related one, so no risk-differentiated review policy, one that
routes high-grade misses to urgent recheck and routine disagreements to
lower-priority audit, can be built on top of it. Increasing accuracy
narrows how often this ambiguity arises without changing what the
output can express, so the limitation is representational rather than a
shortfall to be closed by a better backbone or more data.

This gap is avoidable rather than fundamental, because the structure a
flat output discards is already present in the taxonomy it is trained
on. The Bethesda System defines the atypical squamous cell categories
by changes that are insufficient for a definitive interpretation, and
it separates them by the grade under consideration \cite{nayar2015bethesda}. Read directly,
these two criteria form two axes rather than one flat list, a
grade-related axis and an evidence-definitiveness axis, and the four
abnormal categories are exactly the four combinations of these axes.
Recovering this structure needs no additional annotation beyond the
existing category labels, and it makes expressible a query that a flat
five-class output cannot pose: whether the direct decision is low-grade
while the underlying morphology belief remains high-grade, a
discordance that marks exactly the case a risk-differentiated review
policy needs to catch first.

We build three layers on this factorization. The first trains a
diagnostic model with a direct five-class head alongside
factor-specific supervision, and restricts fusion between the two to
the abnormal conditional distribution, so the boundary between normal
and abnormal cases is preserved by construction rather than by
empirical tuning. The second freezes this diagnostic representation
and maps it through separate adapters into a morphology space and an
evidence space, each factor state described by multiple learnable
prototypes, and treats the resulting separability as something to
audit rather than to assume. The third consumes only the frozen
outputs of the first two layers to add screening-conflict detection,
the discordance sentinel described above, and class-conditional
pairwise prediction sets, so review routing becomes possible without
touching the diagnosis itself. Across five-fold development evaluation
the two axes separate in the direction the underlying clinical
construct predicts, the sentinel recovers the large majority of hidden
high-grade cases at a pre-specified operating point, and the
diagnostic stage matches its unfactorized counterpart under a matched
training budget, so the added structure carries no cost to the
decision it is built on top of.

This paper makes the following contributions.

\begin{itemize}
\item We introduce a factorized diagnostic stage that recodes the four
abnormal Bethesda categories onto a grade-related axis and an
evidence-definitiveness axis read directly from their existing
clinical definitions, and that fuses factor and direct beliefs only on
the abnormal conditional distribution, so the normal-versus-abnormal
screening boundary is preserved by construction rather than by
empirical adjustment.
\item We build separate morphology and evidence semantic spaces over
the frozen diagnostic representation, each factor state represented by
multiple learnable prototypes, and we audit the resulting geometry for
separability and non-collapse rather than assuming it holds.
\item We design a modular risk layer that reads only frozen outputs
from the two stages above to provide a factor-decision discordance
sentinel alongside class-conditional pairwise prediction sets,
enabling review policies that are differentiated by risk type without
altering the underlying diagnosis.
\item We show in five-fold development evaluation that the two
semantic axes separate in the direction predicted by their differing
observability, that the discordance sentinel recovers the large
majority of hidden high-grade cases at a pre-specified operating
point, and that the factorized diagnostic stage matches its
unfactorized counterpart under a matched training budget, indicating
that the added structure comes at no cost to the diagnosis it builds
on.
\end{itemize}

These contributions are embodied in MERA-Dx (Morphology--Evidence Factorized and Risk-Aware Diagnostic Network), which preserves a standard five-class diagnostic output while exposing the factor structure that underlies it and routing risk-differentiated review without touching that output.

\section{Related Work}
\label{sec:related}
\subsection{Deep Learning for Cervical Cytology Classification}
\label{sec:rw_backbone}
Early cervical cytology classification relied on hand-crafted morphometric descriptors, nuclear shape and texture statistics passed to conventional classifiers, and was evaluated on benchmark collections such as the Pap-smear corpus and SIPaKMeD \cite{jantzen2005pap,plissiti2018sipakmed}. Deep representation learning replaced this feature engineering, and the field has since tracked general image recognition closely. Residual and densely connected convolutional networks established the first strong baselines \cite{he2016deep,huang2017densenet}, compound scaling improved the accuracy-to-cost tradeoff \cite{tan2019efficientnet}, and ConvNeXt modernized the convolutional design under a transformer-era training recipe \cite{liu2022convnext}. Vision transformers then introduced global attention over image patches \cite{dosovitsky2021vit}, hierarchical variants such as Swin and MaxViT restored the multi-scale locality that fine-grained cytology depends on \cite{liu2021swin,tu2022maxvit}, and MetaFormer-style analyses indicated that much of the gain follows from the macro architecture rather than the specific token-mixing operator \cite{yu2022metaformer}.

Progress along this line is reported almost exclusively as flat five-class accuracy or macro-averaged F1. Because the label set is consumed as an unordered list of five alternatives, every confusion contributes equally to the reported score, and the trained model retains no internal quantity that separates a grade-related error from an evidence-related one. Stronger backbones reduce how often such confusions occur without changing what the output is able to report about them. Our factorized diagnostic stage addresses this by recoding the four abnormal categories onto two axes taken from their existing clinical definitions, so that grade and evidence beliefs are exposed separately while the direct five-class pathway remains intact. We hold the backbone fixed to Swin-Tiny across every stage of this work so that the reported differences are attributable to this factorization rather than to a change of architecture.

\subsection{Structured Supervision and Label Semantics}
\label{sec:rw_multitask}
Introducing auxiliary targets to shape a shared representation is a long-established strategy. Multi-task learning treats related objectives as mutual regularizers \cite{caruana1997multitask}, and uncertainty-based weighting balances heterogeneous losses without manual tuning of their relative scales \cite{kendall2018multitask}. In medical imaging this typically takes the form of an additional annotation layer, for instance a segmentation mask, a lesion attribute, or a secondary grading scheme, collected alongside the primary diagnostic label.

Such designs improve representation quality but shift the cost to annotation, and the auxiliary target is external to the diagnostic taxonomy rather than derived from it. Consequently, any structure the auxiliary head learns is not guaranteed to correspond to a distinction the clinical reporting standard actually makes. Our factor targets differ in kind: they are deterministic recodings of the existing five-class label, requiring no additional annotation, so a gain from them can only come from re-exposing structure already implicit in the taxonomy. This also fixes the semantics of each axis in advance, which is what later allows a disagreement between the factor beliefs and the direct decision to be read as a specific and clinically meaningful event.

\subsection{Prototype and Metric Learning for Semantic Structure}
\label{sec:rw_prototype}
Metric and prototype methods organize a representation by distance rather than by a linear decision boundary. Prototypical networks summarize each class as a single embedding centroid and classify by proximity \cite{snell2017prototypical}, triplet objectives shape the embedding geometry directly through relative distance constraints \cite{schroff2015facenet}, and supervised contrastive learning extends this to full label supervision with many positives per anchor \cite{khosla2020supervised}. The appeal in a medical setting is that the resulting geometry is inspectable: nearest neighbours of a prototype can be shown to a reader.

These methods, however, typically build a single embedding space shared by all classes and then use it to classify. Two consequences follow. A single space forces one geometry to serve several distinctions at once, so it cannot be asked whether one specific factor is separable independently of another; and because the prototypes drive prediction, any weakness in the geometry propagates directly into diagnostic error. Our semantic stage separates these concerns. We build two factor-specific spaces through independent adapters over a frozen diagnostic representation, each factor state carried by multiple prototypes, and we treat separability as a quantity to audit rather than a mechanism to classify with. Because the diagnosis is frozen, a weak axis is reported as a weak axis instead of degrading the decision.

\subsection{Uncertainty Quantification and Selective Prediction}
\label{sec:rw_uncertainty}
Deciding when not to answer is a separate problem from deciding what to answer. Selective prediction formalizes an abstention option for an otherwise fixed classifier and characterizes the resulting risk-coverage tradeoff \cite{geifman2017selective}, and later work integrates the reject option into training so that the coverage target is optimized jointly with the decision rule \cite{geifman2019selectivenet}. Conformal prediction instead returns a set rather than a single label, with finite-sample coverage guarantees that require only exchangeability rather than a correctly specified model \cite{vovk2005algorithmic,angelopoulos2023conformal}. For classification, set-valued predictors can be constructed to minimize ambiguity at a bounded error level \cite{sadinle2019least}, and adaptive variants improve conditional coverage on hard instances \cite{romano2020classification}. Mondrian conformal prediction restricts calibration to a predefined taxonomy of the sample space \cite{vovk2003mondrian}, and the class-conditional case permits a coverage statement about a specific pair of confusable categories rather than a marginal statement across all classes.

What these methods share is that they operate on a scalar ordering of the model's own confidence. An abstention rule fires when confidence falls below a threshold, and a conformal set widens when no label dominates, but neither construction can distinguish the reason for the ambiguity, since both consume the same flat posterior that produced it. In a setting where one confusion is a missed precursor and another is routine interpretive variation, a single coverage or abstention budget therefore spends the same review effort on both. Our risk layer keeps these guarantees but separates the routes. Class-conditional sets are calibrated within each adjacent pair, while a distinct sentinel monitors disagreement between the grade-related factor belief and the direct decision, a signal that exists only because the two are estimated separately. Both read frozen outputs, so neither revises the diagnostic label they annotate.

\section{Method}
\label{sec:method}

MERA-Dx consists of three stages. The first (Sections~\ref{sec:task}--\ref{sec:factorized}) trains a factorized diagnostic model with a direct five-class head alongside morphology and evidence factor heads, fusing their beliefs only on the abnormal conditional distribution. The second (Section~\ref{sec:semantic}) freezes the diagnostic representation and maps it through separate adapters into factor-specific semantic spaces with multi-prototype structure. The third (Section~\ref{sec:risk}) consumes only frozen outputs to provide screening safety, a factor-decision discordance sentinel, and class-conditional pairwise prediction sets. The three stages are trained sequentially, and the second and third leave the diagnostic output untouched.

\subsection{Grade and Evidence Factor Grid}
\label{sec:task}
Let $x$ denote a single-cell image and let $y$ be one of the five Bethesda
System diagnosis labels. The four abnormal labels are deterministically
mapped onto two factors,
\begin{equation}
    m(y)\in\{\mathrm{Low},\mathrm{High}\}, \qquad
    e(y)\in\{\mathrm{Ambiguous},\mathrm{Definitive}\},
    \label{eq:factors}
\end{equation}
where $m$ is the grade-related morphology factor and $e$ is the
evidence-definitiveness factor. The mapping follows directly from the
Bethesda System's own category definitions \cite{nayar2015bethesda},
\begin{equation}
\begin{array}{c|cc}
    & \mathrm{Ambiguous} & \mathrm{Definitive}\\ \hline
\mathrm{Low} & \mathrm{ASC\!-\!US} & \mathrm{LSIL}\\
\mathrm{High} & \mathrm{ASC\!-\!H} & \mathrm{HSIL}
\end{array}
\label{eq:grid}
\end{equation}
Normal samples remain in the direct five-class pathway but are excluded
from the abnormal factor losses and from the prototype-state updates
described in Section~\ref{sec:semantic}. The factor labels are
deterministic transformations of $y$ motivated by diagnostic semantics,
not additional manual annotations or independently validated biological
categories.

\subsection{Flat Baseline}
\label{sec:baseline}
The flat baseline uses a Swin-Tiny backbone with a single five-class
linear head trained with cross-entropy loss. Images are resized with
fixed letterbox preprocessing to $224\times224$ pixels, and we use a
fixed five-fold development split with seed 42. Content-level duplicate
exclusion is audited, but patient and slide identifiers are not
populated in the supplied manifest, so patient or slide group
disjointness is not claimed. The flat baseline serves only as the
direct diagnostic reference; it is not treated as a semantic or
risk-control module.

\subsection{Morphology-Evidence Factorized Diagnosis}
\label{sec:factorized}
Given backbone features $h=f_{\theta}(x)$, the direct head produces
five-class logits $z_d$. Two factor projectors produce normalized
semantic features,
\begin{equation}
    z_m=A_m(h), \qquad z_e=A_e(h),
    \label{eq:factorproject}
\end{equation}
feeding two-class morphology and evidence heads. The factor losses
provide structured supervision to the shared representation while the
direct five-class loss remains active throughout training.

For the four abnormal classes, the morphology and evidence logits are
combined through a checkerboard interaction term. With $i\in\{0,1\}$
denoting the Low/High morphology state and $j\in\{0,1\}$ denoting the
Ambiguous/Definitive evidence state, the joint abnormal logit is
\begin{equation}
    \ell_{ij}=\ell^m_i+\ell^e_j+\eta\,s_{ij}, \qquad
    s=\begin{bmatrix}+1&-1\\-1&+1\end{bmatrix}.
    \label{eq:interaction}
\end{equation}
The scalar interaction $\eta$ is produced from the factor features,
bounded by a hyperbolic tangent, and initialized at zero. This
checkerboard parameterization removes the common-shift degree of
freedom and isolates a single interaction contrast across the four
abnormal combinations, since a common shift applied to all four logits
cancels under softmax while the checkerboard term changes their
relative combination.

The final fusion acts only on the abnormal conditional distribution.
Let $p_0$ denote the Normal probability from the direct head,
$p_A=1-p_0$ the abnormal mass, $q_d$ the direct abnormal conditional
distribution, and $q_T$ the factorized distribution from
Eq.~\eqref{eq:interaction}. The fused conditional distribution is a
geometric interpolation,
\begin{equation}
    q_f=\operatorname{softmax}\left((1-\alpha)\log q_d+\alpha\log q_T\right),
    \qquad 0<\alpha<0.30,
    \label{eq:fusion}
\end{equation}
and the final five-class probabilities are $p_f=[p_0,\, p_Aq_f]$. The
Normal/Abnormal mass is therefore preserved exactly by construction,
not by empirical tuning. We audit the abnormal-row sum, the
screening-mass error, the interaction magnitude, and $\alpha$ in every
fold.

\subsection{Dual-Space Semantic Geometry}
\label{sec:semantic}
The factorized diagnosis stage above is then frozen in its entirety,
including its direct probabilities and factor heads. The frozen
feature $h$ is mapped into two separate factor-specific spaces,
\begin{equation}
    u_m=\operatorname{norm}(g_m(h)), \qquad
    u_e=\operatorname{norm}(g_e(h)),
    \label{eq:spaces}
\end{equation}
where $g_m$ and $g_e$ are separate adapters, one per factor. Each
factor state is represented by two learnable prototypes rather than a
single point, leaving room for within-state variation. For factor
$r\in\{m,e\}$ and state $s\in\{0,1\}$, the state score is a
temperature-scaled log-sum-exp over its two prototypes,
\begin{equation}
    a^r_s(u_r)=\tau\log\sum_{k=1}^{2}
    \exp\left(\frac{\operatorname{cos}(u_r,p^r_{s,k})}{\tau}\right),
    \qquad \tau=0.20.
    \label{eq:prototype}
\end{equation}
We do not claim that the two prototypes assigned to a state correspond
to two distinct biological phenotypes; the second prototype is present
to absorb within-state variation, not to name a subtype. A secondary
compositional readout can combine morphology and evidence state scores
into a five-class posterior, but this readout is never fused back into
the diagnosis. This stage therefore answers a representation question
rather than a diagnostic one: do the two factor-specific spaces carry
stable positive separation and non-collapsed prototypes across folds,
independent of the classifier that produced the frozen features.

We report the fold-level morphology and evidence geometry gaps
together with their bootstrap uncertainty intervals, factor-specific
linear-probe contrasts, prototype non-collapse, and confirmation that
the frozen diagnosis is unmodified by this stage. For a factor $r$,
the geometry gap is the mean cosine similarity of pairs sharing state
$r$ minus the mean cosine similarity of pairs differing in state $r$
while sharing the other factor, computed over the two abnormal-class
pairings that isolate that factor. Each term samples at most 50,000
pairs and obtains a 95\% percentile interval from 2,000 bootstrap
resamples of the sampled similarities. The factor-specificity
contrasts are $D_m=A_{m\to m}-A_{m\to e}$ and
$D_e=A_{e\to e}-A_{e\to m}$, where $A_{r\to t}$ is three-fold
cross-validated linear-probe accuracy for target factor $t$ read from
space $r$. The primary evidence at this stage is the geometric
organization itself, not the accuracy of the secondary prototype
readout.

\subsection{Modular Risk Layer}
\label{sec:risk}
This layer consumes only frozen outputs from the factorized diagnosis
stage above. It does not require the semantic-geometry readout, does
not train a new diagnostic classifier, and does not change the frozen
top-one label, probabilities, logits, or parameters. Its output
augments the frozen diagnosis with a prediction set, screening and
high-grade probabilities, uncertainty fields, risk flags, a review
priority, an abstention decision, and a reason code.

\textbf{Screening safety.} This component applies the two-state
Normal/Abnormal distribution induced by the frozen direct
probabilities. With the class-conditional nonconformity score
$1-q_y$, a binary prediction set is calibrated at nominal coverage
$\gamma_{\mathrm{scr}}=0.90$ on the inner calibration pool. A
screening-boundary flag is raised when the set contains a state
opposing the frozen top-one prediction; the development gate requires
screening coverage of at least 0.90 and a raw abnormal-to-normal
error rate of at most 0.01. Calibration and evaluation rows are
disjoint within each fold.

\textbf{Factor-decision discordance sentinel.} Among direct
predictions outside the high-grade pair, this sentinel identifies
cases where the direct decision is low-grade but the factor beliefs
retain high-grade evidence. The sentinel feature vector is
\begin{equation}
    \phi_{\mathrm{HU}}(x)=
    [p_A,p_H^d,p_H^m,p_{\mathrm{Def}}^e,\Delta_H,M_{\mathrm{severity}},H_d,\eta],
    \label{eq:sentinel_features}
\end{equation}
and its sentinel score is
\begin{equation}
    r_{\mathrm{HU}}(x)=\sigma(w^\top\phi_{\mathrm{HU}}(x)+b).
    \label{eq:sentinel_score}
\end{equation}
Here $p_d$ denotes the five-class probability vector from the frozen
direct head. Thus $p_A=1-p_d(\mathrm{Normal})$ and
$p_H^d=p_d(\mathrm{ASC\!-\!H})+p_d(\mathrm{HSIL})$, while $p_H^m$ is
the morphology-high probability and $p_{\mathrm{Def}}^e$ is the
evidence-definitive probability, both read from the factorized
diagnosis stage. The remaining features are the logit-scale
disagreement $\Delta_H$, the direct high-versus-low margin
$M_{\mathrm{severity}}$, the normalized five-class entropy of the
direct head $H_d$, and the bounded interaction term $\eta$. The
eight-dimensional vector is standardized and fitted with a
class-balanced logistic sentinel using cross-fitted development
predictions, and its operating threshold is fixed by a finite-sample
95\% positive-score calibration rule. The sentinel raises a
high-priority review flag while leaving the diagnosis itself
unchanged. Its target is a hidden high-grade case,
\begin{equation}
 y_{\mathrm{HU}}=\mathbf{1}\!\left[y\in\{\mathrm{ASC\!-\!H},\mathrm{HSIL}\}\;\land\;\hat y\notin\{\mathrm{ASC\!-\!H},\mathrm{HSIL}\}\right],
 \label{eq:hu_target}
\end{equation}
namely a case whose true category is high-grade while the direct
prediction $\hat y$ is not.

\textbf{Class-conditional pairwise sets.} For the low pair
$L=\{\mathrm{ASC\!-\!US},\mathrm{LSIL}\}$, the conditional probability
and nonconformity score are
\begin{equation}
    q_L(c\mid x)=\frac{p_c(x)}{p_{\mathrm{ASC\!-\!US}}(x)+p_{\mathrm{LSIL}}(x)},
    \qquad s(x,c)=1-q_L(c\mid x),
    \label{eq:pair_probability}
\end{equation}
and the high pair is defined analogously. For calibration class $c$
with $n_c$ samples, the nominal calibration target is
$\gamma_{\mathrm{cal}}=0.95$ with finite-sample rank
$k_c=\lceil(n_c+1)\gamma_{\mathrm{cal}}\rceil$. The prediction set is
\begin{equation}
    \Gamma(x)=\{c:s(x,c)\leq Q_c\},
    \label{eq:conformal_set}
\end{equation}
where $Q_c$ is the class-conditional calibration quantile from the
outer-training calibration pool,
\begin{equation}
 Q_c=\begin{cases}s_{(k_c)},& k_c\le n_c,\\1,& k_c>n_c,\end{cases}
 \label{eq:pair_quantile}
\end{equation}
and the development evaluation gate requires evaluation coverage of at
least 0.90. If a raw binary set is empty, the implementation includes
the higher-probability class as a deterministic finite-sample fallback,
and raw empty sets are counted before this fallback is applied.
Calibration and evaluation rows are kept disjoint within each fold. A
genuinely difficult adjacent case therefore receives a two-element set
rather than an unjustified singleton, which is the intended
communication behavior rather than a failure to decide. Because no
untouched externally labeled cohort was available, the evidence from
this layer is exploratory development evidence. This construction is a
class-conditional, or Mondrian, conformal predictor restricted to the
adjacent pair \cite{vovk2003mondrian,sadinle2019least}.

\begin{figure*}[!t]
    \centering
    \includegraphics[width=0.75\textwidth,height=0.55\textheight,keepaspectratio]{figures/overall_framework.png}
    \caption{Overall framework. A single-cell image is encoded by a shared Swin-Tiny backbone into a direct five-class head and two factor heads that read the grade-related and evidence-definitiveness axes recovered from the Bethesda System's category definitions. The two factor beliefs are combined through a bounded checkerboard interaction and fused with the direct abnormal conditional distribution while leaving Normal and Abnormal mass exactly preserved. The frozen representation is then mapped by separate adapters into a morphology space and an evidence space carrying multiple learnable prototypes per factor state, and a modular risk layer consumes only these frozen outputs to provide screening safety, a factor-decision discordance sentinel, and class-conditional pairwise prediction sets.}
    \label{fig:framework}
\end{figure*}

\section{Experimental Protocol}
\label{sec:protocol}

\subsection{Data Partition and Reproducibility}

\begin{table}[t]
\caption{Image-level composition of the cleaned XUData TBS5 cohort.
Counts are single-cell images. The test split was sealed during model
development.}
\label{tab:dataset_composition}
\centering
\footnotesize
\setlength{\tabcolsep}{3.5pt}
\renewcommand{\arraystretch}{0.95}
\begin{tabular}{lrrrrrr}
\toprule
Split & Normal & ASC-US & LSIL & ASC-H & HSIL & Total \\
\midrule
Train        & 1,628 & 1,657 & 1,627 & 1,458 & 1,615 & 7,985 \\
Development  &   203 &   206 &   204 &   183 &   202 &   998 \\
Calibration  &   203 &   206 &   203 &   182 &   202 &   996 \\
Sealed test  &   381 &   389 &   382 &   345 &   380 & 1,877 \\
\midrule
All images   & 2,415 & 2,458 & 2,416 & 2,168 & 2,399 & 11,856 \\
\bottomrule
\end{tabular}
\end{table}

The cleaned XUData TBS5 cohort contains 11,856 single-cell images
(224 $\times$ 224 $\times$ 3). The training split contains 7,985 images,
whereas the development, calibration and sealed-test splits contain 998,
996 and 1,877 images, respectively. The morphology and evidence targets
are deterministically derived from the five-class diagnosis labels; they
are model-defined semantic factors rather than independently pathologist-annotated ground truth.

The locked five-fold OOF analysis used a 7,586-image S0-R pool formed
from the train-fit and select-s0 subsets of the training split.
Therefore, the 7,586 OOF rows are a downstream evaluation subset of the
7,985 training images, rather than an additional dataset split. All five
held-out folds covered the 7,586 S0-R samples exactly once. Patient and
slide identifiers are not populated in the supplied manifest, and
filename-derived candidate groups cross split boundaries, so patient or
slide group-disjointness is not claimed. The
flat baseline, the factorized diagnosis stage, and the semantic-geometry
stage use the same image size and letterbox preprocessing. Aggregate
inventory metadata were available for dataset accounting, but no test
images, predictions, or sample-level test labels were used during model
development or evaluation. Reported values are copied from the
pre-specified result files and fold manifests.

The experiment chain is intentionally staged. The flat baseline records
the reference. The factorized diagnosis stage is trained with the direct
five-class loss and factor losses, and epoch-30 checkpoints are retained
for all folds. The semantic-geometry stage is trained only through
adapters and prototypes while the factorized diagnosis is frozen. The
risk layer consumes frozen outputs and fixed calibration artifacts;
thresholds and epochs are not changed post hoc.

\subsection{Implementation Details}
\label{sec:implementation}
The backbone is ImageNet-initialized Swin-Tiny with $224\times224$
letterbox inputs. MERA-Dx uses batch size 64, AdamW with learning rate
$10^{-4}$, weight decay $10^{-4}$, automatic mixed precision, and seed 42.
The factorized diagnosis stage uses $\lambda_m=0.15$, $\lambda_e=0.15$,
$\lambda_{\mathrm{joint}}=0.10$, and $\lambda_{\eta}=0.01$; its
interaction dimension is 32 and semantic dimension is 128. The bounded
fusion uses $\alpha_{\max}=0.30$ and initialization $\alpha_0=0.05$. The
semantic-geometry stage uses batch size 32, 12 epochs, adapter learning
rate $5\times10^{-5}$, prototype temperature $0.20$, two prototypes per
factor state, and a frozen factorized-diagnosis checkpoint. The risk
layer adds no large image classifier and uses only frozen outputs plus
fixed calibration artifacts.

The competing methods listed in Table~\ref{tab:backbone_comparison} were
retrained on XUData TBS5 following their published architectures and training
recipes, under the same five-fold partition, preprocessing, random seed, and
five-class label mapping used for MERA-Dx.
The supplied record does not contain verified acquisition, stain, resolution,
hardware, or ethics/access metadata; these fields are not asserted in the
present manuscript and remain a prerequisite for a complete submission.

\subsection{Metrics}
For diagnosis, we report accuracy, macro-F1, abnormal macro-F1, low-pair
and high-pair binary macro-F1, and screening sensitivity. Pair metrics
are computed after renormalizing the corresponding two-class probability
row. For the semantic-geometry audit, we report the cosine-similarity
geometry gaps and their deterministic 2,000-resample 95\% percentile
intervals, factor-specific three-fold linear-probe contrasts, and
parameter drift. For the risk layer, we report screening coverage,
high-grade undercall review recall, low and high pair coverage, raw
empty-set rate, review-priority rates, and the number of prediction
changes.

\subsection{Evaluation Boundary}
The flat baseline, factorized diagnosis, and semantic-geometry results
are five-fold development results. The risk layer is explicitly marked
exploratory and is not a formal external confirmation. No test images,
predictions, or sample-level test labels were used in model development
or evaluation, and no untouched external TBS5 confirmation cohort was
available.

\section{Results}
\label{sec:results}

\subsection{Comparison on XUData TBS5}
\label{sec:backbone_comparison}
Table~\ref{tab:backbone_comparison} compares MERA-Dx against nine published
cervical cytology classification methods, all retrained and evaluated on the
XUData TBS5 development cohort under the same five-fold protocol,
preprocessing, and label mapping. MERA-Dx reaches accuracy 0.7776 and
macro-F1 0.7752, ahead of the strongest competing method A2SDNet121
(0.7080 / 0.7068) by 6.96 and 6.84 percentage points respectively. The
remaining methods fall between 0.5588 and 0.6846 macro-F1, with the
handcrafted feature-selection pipelines (PCA+GWO, CNN+GA+SVM) trailing the
end-to-end deep models. No test or external confirmation data were used.

\begin{table*}[!t]
\centering
\small
\caption{Comparison with existing cervical cytology classification methods on
the XUData TBS5 development cohort. All methods were retrained on XUData
under the same five-fold protocol, preprocessing, and five-class label
mapping. Values are means across the five fixed development folds.
Patient/slide group-disjointness was not verified.}
\label{tab:backbone_comparison}
\setlength{\tabcolsep}{4pt}
\begin{tabular*}{\textwidth}{@{\extracolsep{\fill}}lccccccc@{}}
\toprule
\textbf{Method} & \textbf{Acc} & \textbf{Bal-Acc} & \textbf{Macro-Prec} & \textbf{Macro-Rec} & \textbf{Macro-Spec} & \textbf{Macro-F1} & \textbf{Macro-AUC} \\
\midrule
HCT-Net         & 0.6268 & 0.6257 & 0.6427 & 0.6257 & 0.9067 & 0.6291 & 0.8945 \\
LGPNet          & 0.6784 & 0.6767 & 0.6870 & 0.6767 & 0.9196 & 0.6778 & 0.9174 \\
A2SDNet121      & 0.7080 & 0.7061 & 0.7121 & 0.7061 & 0.9271 & 0.7068 & 0.9276 \\
DeepCervix-HDFF & 0.6835 & 0.6818 & 0.6916 & 0.6818 & 0.9209 & 0.6820 & 0.9173 \\
MSCCNet         & 0.6879 & 0.6861 & 0.6907 & 0.6861 & 0.9220 & 0.6846 & 0.9167 \\
CerCan-Net      & 0.6743 & 0.6724 & 0.6796 & 0.6724 & 0.9186 & 0.6713 & 0.9094 \\
DIFF            & 0.6524 & 0.6509 & 0.6693 & 0.6509 & 0.9131 & 0.6545 & 0.9064 \\
PCA+GWO         & 0.5718 & 0.5728 & 0.5658 & 0.5728 & 0.8930 & 0.5588 & 0.8421 \\
CNN+GA+SVM      & 0.6627 & 0.6616 & 0.6610 & 0.6616 & 0.9156 & 0.6605 & 0.9046 \\
\midrule
\textbf{MERA-Dx (ours)} & \textbf{0.7776} & \textbf{0.7758} & \textbf{0.7760} & \textbf{0.7758} & \textbf{0.9445} & \textbf{0.7752} & \textbf{0.9459} \\
\bottomrule
\end{tabular*}
\end{table*}

\subsection{Comparison on the Public Benchmark}
\label{sec:public_comparison}
Table~\ref{tab:public_comparison} repeats the comparison on a public cervical
cytology benchmark under the same five-fold protocol and label mapping.
MERA-Dx reaches accuracy 0.9539 and macro-F1 0.9533, ahead of the strongest
competing method DeepCervix-HDFF (0.9489 / 0.9476), with the highest MCC
(0.9420) and macro-AUROC (0.9960) in the comparison. The ordering of the
competing methods differs from the XUData result, which is consistent with the
public benchmark being a substantially easier and more class-balanced
collection.

\begin{table*}[!t]
\centering
\small
\caption{Comparison with existing cervical cytology classification methods on
the public benchmark. All methods were retrained under the same five-fold
protocol, preprocessing, and label mapping. Values are means across the five
folds.}
\label{tab:public_comparison}
\setlength{\tabcolsep}{4pt}
\begin{tabular*}{\textwidth}{@{\extracolsep{\fill}}lcccccccc@{}}
\toprule
\textbf{Method} & \textbf{Acc} & \textbf{Bal-Acc} & \textbf{Macro-Prec} & \textbf{Macro-Rec} & \textbf{Macro-Spec} & \textbf{Macro-F1} & \textbf{MCC} & \textbf{Macro-AUC} \\
\midrule
HCT-Net         & 0.8476 & 0.8489 & 0.8479 & 0.8489 & 0.9620 & 0.8434 & 0.8110 & 0.9746 \\
LGPNet          & 0.9313 & 0.9309 & 0.9299 & 0.9309 & 0.9829 & 0.9296 & 0.9140 & 0.9909 \\
A2SDNet121      & 0.8223 & 0.8250 & 0.8228 & 0.8250 & 0.9556 & 0.8177 & 0.7800 & 0.9580 \\
DeepCervix-HDFF & 0.9489 & 0.9488 & 0.9474 & 0.9488 & 0.9873 & 0.9476 & 0.9360 & 0.9951 \\
MSCCNet         & 0.9392 & 0.9399 & 0.9376 & 0.9399 & 0.9849 & 0.9379 & 0.9240 & 0.9942 \\
CerCan-Net      & 0.9321 & 0.9336 & 0.9309 & 0.9336 & 0.9831 & 0.9313 & 0.9150 & 0.9916 \\
DIFF            & 0.9201 & 0.9204 & 0.9179 & 0.9204 & 0.9801 & 0.9178 & 0.9000 & 0.9851 \\
PCA+GWO         & 0.9175 & 0.9194 & 0.9147 & 0.9194 & 0.9794 & 0.9152 & 0.8970 & 0.9828 \\
CNN+GA+SVM      & 0.9266 & 0.9263 & 0.9244 & 0.9263 & 0.9817 & 0.9248 & 0.9080 & 0.9905 \\
\midrule
\textbf{MERA-Dx (ours)} & \textbf{0.9539} & \textbf{0.9552} & \textbf{0.9528} & \textbf{0.9552} & \textbf{0.9885} & \textbf{0.9533} & \textbf{0.9420} & \textbf{0.9960} \\
\bottomrule
\end{tabular*}
\end{table*}

\subsection{Diagnostic-Path Comparison}
\label{sec:path_comparison}
MERA-Dx's diagnostic stage provides a direct-versus-final path
comparison; the values are nearly identical. It does not isolate
individual auxiliary losses.

\begin{table}[!htbp]
\caption{Diagnostic-path comparison. Y indicates an active component; the
direct row is evaluated before the final mass-preserving fusion.}
\label{tab:path}
\centering
\footnotesize
\setlength{\tabcolsep}{2.3pt}
\begin{tabular}{lcccccc}
\toprule
Variant & M & E & $\eta$ & Fusion & Macro-F1 & Low / High \\
\midrule
Flat baseline$^\dagger$ & -- & -- & -- & -- & 0.7576 & 0.7889 / 0.7726 \\
Factorized direct & Y & Y & Y & -- & 0.7748 & 0.7979 / 0.7953 \\
Factorized final & Y & Y & Y & Y & 0.7752 & 0.7979 / 0.7953 \\
\bottomrule
\end{tabular}
\vskip 2pt
{\scriptsize $\dagger$ Historical epoch-12 reference.}
\end{table}

\subsection{Semantic Geometry Shows Positive Factor-Specific Separation}
Table~\ref{tab:geometry} reports the primary semantic-geometry result.
The morphology geometry gap is $1.2627\pm0.0664$ and the evidence gap is
$0.5184\pm0.0528$. The minimum fold-level lower bounds exported by the
audit are 1.1735 and 0.4667, respectively, both strictly positive. The
factor-specificity contrasts are $D_m=0.3625$ and $D_e=0.1694$ on
average, and all five folds pass the geometry gate. These values are
descriptive outputs of the five-fold development audit; they should not
be read as independent biological validation of the recoded factors.

\begin{table}[!htbp]
\caption{Dual-space semantic geometry audit.}
\label{tab:geometry}
\centering
\footnotesize
\setlength{\tabcolsep}{3.0pt}
\begin{tabular}{lc}
\toprule
Audit quantity & Five-fold result \\
\midrule
Morphology gap & $1.2627\pm0.0664$ \\
Evidence gap & $0.5184\pm0.0528$ \\
Minimum morphology interval lower bound & 1.1735 \\
Minimum evidence interval lower bound & 0.4667 \\
$D_m$ & 0.3625 \\
$D_e$ & 0.1694 \\
Geometry gate & $5/5$ folds pass \\
Parameter drift & $0/5$ folds \\
Prediction changes & $0/5$ folds \\
\bottomrule
\end{tabular}
\end{table}

The prototype readout is weaker than the factorized diagnosis
(Macro-F1 approximately 0.7101), so the semantic-geometry stage is an
analysis layer rather than a replacement classifier. It organizes the
frozen representation into factor-specific spaces with positive gaps,
non-collapsed prototypes, and no diagnostic change.

\subsection{Risk Layer Provides Exploratory Decision Support}
The risk-layer modules target distinct risks rather than reclassifying
every sample. As shown in Table~\ref{tab:risk}, screening coverage is
99.62\% on average, high-grade undercall review recall is 94.70\% with a
worst-fold value of 93.33\%, and low-pair and high-pair prediction-set
coverage is 95.10\% and 94.86\%. Across the five folds, 266 high-undercall
events were identified and 252 were captured by the sentinel review flag.
The minimum class-conditional pair coverage across ASC-US, LSIL, ASC-H,
and HSIL was 93.0\%. Here undercall recall is defined over true ASC-H or
HSIL samples whose frozen top-one prediction is outside the high pair;
pair coverage is aggregated over the corresponding two-class subsets. The
raw empty-set rate is zero. Parameter drift and prediction changes are
both zero, confirming that the safety layer is non-invasive.

\begin{table}[!htbp]
\caption{Exploratory risk outputs over five folds.}
\label{tab:risk}
\centering
\footnotesize
\setlength{\tabcolsep}{3.0pt}
\begin{tabular}{lc}
\toprule
Risk quantity & Mean result \\
\midrule
Screening coverage & 99.62\% \\
High-undercall review recall & 94.70\% \\
High-undercall events captured & 252/266 \\
Worst-fold high-undercall recall & 93.33\% \\
Low-pair coverage & 95.10\% \\
High-pair coverage & 94.86\% \\
Minimum class-conditional pair coverage & 93.0\% \\
Raw abnormal-to-normal top-1 error & 0.28\% \\
Raw empty-set rate & 0\% \\
Silent high undercall & 0.48\% \\
Parameter drift & 0 \\
Prediction changes & 0 \\
\bottomrule
\end{tabular}
\end{table}

The 43.41\% overall review rate is not a severe-failure rate: the router
separates safety-critical abstention (18.09\%) from medium-priority
adjacent ambiguity, while preserving the frozen top-one label.

No matched-review-budget comparison with maximum-softmax, entropy, or
direct high-grade scores is present in the available result files;
therefore, the risk layer is not claimed to outperform those baselines.

\section{Discussion and Limitations}
\label{sec:discussion}
The results link three levels: the factorized diagnosis improves the
diagnostic performance while preserving screening mass; the
semantic-geometry audit yields positive dual-space separation without
changing the diagnosis; and the risk layer exposes factor discordance and
conditional prediction sets.

Claims are limited by the development-only five-fold design, absent
patient and slide identifiers, filename-derived groups crossing splits,
and unavailable acquisition, preparation, stain, hardware,
label-adjudication, and IRB or consent metadata. No patient-level
generalization or independent test performance is claimed. The
semantic-geometry stage is representation analysis and the risk layer
exploratory development behavior; screening and undercall quantities are
cell-level, and factor labels require prospective pathologist validation.

The next scientifically meaningful validation is not another threshold
search on the current results. It is a preregistered evaluation on a new
patient or slide group-disjoint TBS5 cohort or a temporal holdout that
has not been exposed to model development, calibration, smoke tests, or
model selection. Until that resource is available, the appropriate claim
is exploratory risk-aware decision support.

\section{Conclusion}
\label{sec:conclusion}
We presented a staged framework linking factorized diagnosis, dual-space
semantic auditing, and risk-aware decision support. The factorized
diagnosis improves matched-budget performance without changing screening
mass; the semantic-geometry audit reports factor separation without
modifying the diagnosis; and the risk layer reports screening,
high-grade, and adjacent-category risks while preserving the frozen
diagnosis. Findings remain exploratory pending external confirmation.

\section*{Ethics and Data Availability Note}
This study uses the XUData TBS5 cohort. The supplied experiment artifacts
do not contain institution-specific ethics identifiers, patient-consent
language, or data-access restrictions; accordingly, no unverified ethics
or consent statement is made here. The reported analyses use development
out-of-fold rows only; test images, predictions, and sample-level test
labels were not used. Institution-specific ethics and data-access wording
must be verified against the approved dataset documentation before
submission.

\FloatBarrier
\begin{thebibliography}{00}
\footnotesize
\setlength{\itemsep}{0pt}
\setlength{\parsep}{0pt}

\bibitem{nayar2015bethesda}
R. Nayar and D. C. Wilbur, Eds., \emph{The Bethesda System for Reporting Cervical Cytology: Definitions, Criteria, and Explanatory Notes}, 3rd ed. Cham, Switzerland: Springer, 2015.

\bibitem{jantzen2005pap}
J. Jantzen, J. Norup, G. Dounias, and B. Bjerregaard, ``Pap-smear benchmark data for pattern classification,'' in \emph{Proc. Nature Inspired Smart Inf. Syst. (NiSIS)}, 2005, pp. 1--9.

\bibitem{plissiti2018sipakmed}
M. E. Plissiti, P. Dimitrakopoulos, G. Sfikas, C. Nikou, O. Krikoni, and A. Charchanti, ``SIPaKMeD: A new dataset for feature and image based classification of normal and pathological cervical cells in Pap smear images,'' in \emph{Proc. IEEE Int. Conf. Image Process. (ICIP)}, 2018, pp. 3144--3148.

\bibitem{he2016deep}
K. He, X. Zhang, S. Ren, and J. Sun, ``Deep residual learning for image recognition,'' in \emph{Proc. IEEE Conf. Comput. Vis. Pattern Recognit. (CVPR)}, 2016, pp. 770--778.

\bibitem{huang2017densenet}
G. Huang, Z. Liu, L. van der Maaten, and K. Q. Weinberger, ``Densely connected convolutional networks,'' in \emph{Proc. IEEE Conf. Comput. Vis. Pattern Recognit. (CVPR)}, 2017, pp. 4700--4708.

\bibitem{tan2019efficientnet}
M. Tan and Q. V. Le, ``EfficientNet: Rethinking model scaling for convolutional neural networks,'' in \emph{Proc. 36th Int. Conf. Mach. Learn. (ICML)}, vol. 97, 2019, pp. 6105--6114.

\bibitem{liu2022convnext}
Z. Liu, H. Mao, C.-Y. Wu, C. Feichtenhofer, T. Darrell, and S. Xie, ``A ConvNet for the 2020s,'' in \emph{Proc. IEEE/CVF Conf. Comput. Vis. Pattern Recognit. (CVPR)}, 2022, pp. 11976--11986.

\bibitem{dosovitsky2021vit}
A. Dosovitskiy \emph{et al.}, ``An image is worth 16x16 words: Transformers for image recognition at scale,'' in \emph{Proc. Int. Conf. Learn. Represent. (ICLR)}, 2021.

\bibitem{liu2021swin}
Z. Liu, Y. Lin, Y. Cao, H. Hu, Y. Wei, Z. Zhang, S. Lin, and B. Guo, ``Swin Transformer: Hierarchical vision transformer using shifted windows,'' in \emph{Proc. IEEE/CVF Int. Conf. Comput. Vis. (ICCV)}, 2021, pp. 10012--10022.

\bibitem{tu2022maxvit}
Z. Tu, H. Talebi, H. Zhang, F. Yang, P. Milanfar, A. Bovik, and Y. Li, ``MaxViT: Multi-axis vision transformer,'' in \emph{Proc. Eur. Conf. Comput. Vis. (ECCV)}, 2022, pp. 459--479.

\bibitem{yu2022metaformer}
W. Yu, C. Si, P. Zhou, M. Luo, Y. Zhou, J. Feng, S. Yan, and X. Wang, ``MetaFormer baselines for vision,'' arXiv preprint arXiv:2210.13452, 2022.

\bibitem{caruana1997multitask}
R. Caruana, ``Multitask learning,'' \emph{Mach. Learn.}, vol. 28, no. 1, pp. 41--75, 1997.

\bibitem{kendall2018multitask}
A. Kendall, Y. Gal, and R. Cipolla, ``Multi-task learning using uncertainty to weigh losses for scene geometry and semantics,'' in \emph{Proc. IEEE/CVF Conf. Comput. Vis. Pattern Recognit. (CVPR)}, 2018, pp. 7482--7491.

\bibitem{snell2017prototypical}
J. Snell, K. Swersky, and R. S. Zemel, ``Prototypical networks for few-shot learning,'' in \emph{Adv. Neural Inf. Process. Syst. (NeurIPS)}, vol. 30, 2017, pp. 4077--4087.

\bibitem{schroff2015facenet}
F. Schroff, D. Kalenichenko, and J. Philbin, ``FaceNet: A unified embedding for face recognition and clustering,'' in \emph{Proc. IEEE Conf. Comput. Vis. Pattern Recognit. (CVPR)}, 2015, pp. 815--823.

\bibitem{khosla2020supervised}
P. Khosla, P. Teterwak, C. Wang, A. Sarna, Y. Tian, P. Isola, A. Maschinot, C. Liu, and D. Krishnan, ``Supervised contrastive learning,'' in \emph{Adv. Neural Inf. Process. Syst. (NeurIPS)}, vol. 33, 2020, pp. 18661--18673.

\bibitem{geifman2017selective}
Y. Geifman and R. El-Yaniv, ``Selective classification for deep neural networks,'' in \emph{Adv. Neural Inf. Process. Syst. (NeurIPS)}, vol. 30, 2017, pp. 4878--4887.

\bibitem{geifman2019selectivenet}
Y. Geifman and R. El-Yaniv, ``SelectiveNet: A deep neural network with an integrated reject option,'' in \emph{Proc. 36th Int. Conf. Mach. Learn. (ICML)}, vol. 97, 2019, pp. 2151--2159.

\bibitem{vovk2005algorithmic}
V. Vovk, A. Gammerman, and G. Shafer, \emph{Algorithmic Learning in a Random World}. New York, NY, USA: Springer, 2005.

\bibitem{vovk2003mondrian}
V. Vovk, D. Lindsay, I. Nouretdinov, and A. Gammerman, ``Mondrian confidence machine,'' Royal Holloway, Univ. London, Tech. Rep., 2003.

\bibitem{sadinle2019least}
M. Sadinle, J. Lei, and L. Wasserman, ``Least ambiguous set-valued classifiers with bounded error levels,'' \emph{J. Amer. Statist. Assoc.}, vol. 114, no. 525, pp. 223--234, 2019.

\bibitem{romano2020classification}
Y. Romano, M. Sesia, and E. J. Cand\`{e}s, ``Classification with valid and adaptive coverage,'' in \emph{Adv. Neural Inf. Process. Syst. (NeurIPS)}, vol. 33, 2020, pp. 3581--3591.

\bibitem{angelopoulos2023conformal}
A. N. Angelopoulos and S. Bates, ``Conformal prediction: A gentle introduction,'' \emph{Found. Trends Mach. Learn.}, vol. 16, no. 4, pp. 494--591, 2023.

\end{thebibliography}

\end{document}
